% P0030 Yvyra Soy — LaTeX submission template
%
% Compile: latexmk -pdf paper.tex
%
% Target journal: Computers and Electronics in Agriculture (Elsevier)
% Subject: Indigenous territory × soybean frontier in Paraguay
%
% Data lineage (all real, public):
%   - INE 2022 census (557 communities, 19 pueblos)
%   - OSM indigenous community polygons (30 features, admin_level 10/11)
%   - MAG cereal yield 2007/08-2024/25 (soybean by department)
%   - Hansen GFC v1.11 (off-repo, ~56 km^2 tile used as anchor)
%
% Data integrity note (2026-09-09): The thesis author has not contacted
% any indigenous community; this is a public-data analysis only.
% No FPIC engagement has been conducted. See ACTUAL_RESULTS.md for
% data provenance and the audit 15-hat document for ethics context.

\documentclass[preprint,linenumbers]{elsarticle}

\usepackage[utf8]{inputenc}
\usepackage[english, spanish]{babel}
\usepackage{amsmath, amssymb}
\usepackage{graphicx}
\usepackage{booktabs}
\usepackage{siunitx}
\usepackage[colorlinks=true, linkcolor=blue, citecolor=blue, urlcolor=blue]{hyperref}

\bibliographystyle{elsarticle-num}

\begin{document}

\begin{frontmatter}

\title{Yvyra Soy: Indigenous territory presence and soybean frontier
expansion in Paraguay's Gran Chaco and Eastern regions --- \\nA public-data
analysis using INE 2022 census, OSM community polygons, and MAG
department yield (2007/08--2024/25)}

\author[una]{Iv\'an Weiss Van der Pol}
\affiliation[una]{organization={FP-UNA (Facultad Polit\'ecnica), Universidad Nacional de Asunci\'on},addressline={Asunci\'on},country={Paraguay}}

\begin{abstract}
We present \textbf{Yvyra Soy}, a public-data analysis of the
relationship between indigenous territory presence and soybean
frontier expansion across Paraguay's 16 soy-producing departments,
2007/08--2024/25. The pipeline joins three openly-published data
sources: the \textbf{INE 2022 indigenous communities census}
(557 communities across 19 pueblos), \textbf{OSM community polygons}
(30 features with admin\_level 10/11), and \textbf{MAG cereal yield}
records for soybean (\textit{Glycine max}) by department-year.

\textbf{Measured pilot results (see \texttt{ACTUAL\_RESULTS.md}):}
\begin{itemize}
\item \textbf{Top pueblos by community count:} Mbya Guaran\'i (200),
Ava Guaran\'i (150), Pai Tavytera (61), Ayoreo (26), Nivacle (23).
\item \textbf{National mean soy yield (5-year):} \SI{2429}{kg/ha}
(weighted by department area; range \SI{1900}{kg/ha}--\SI{3100}{kg/ha}).
\item \textbf{Top departments by yield:} Amambay (\SI{3095}{kg/ha}),
Alto Paraguay (\SI{2982}{kg/ha}), Paraguar\'i (\SI{2930}{kg/ha}),
Cordillera (\SI{2742}{kg/ha}), Guair\'a (\SI{2662}{kg/ha}).
\end{itemize}

We \textbf{do not} claim a causal link between indigenous territory
presence and soybean expansion. The headline contribution is
\emph{infrastructure}: a reproducible, public-data pipeline that
joins administrative census records with agricultural statistics
and OpenStreetMap community polygons, all under MIT license.
This infrastructure can be reused by Paraguayan agencies (INFONA,
INDI, MAG) for evidence-based monitoring of the agricultural
frontier in indigenous lands, in compliance with CARE Principles
once FPIC engagement is established (currently zero per the 2026-09-09
audit; see \texttt{ACTUAL\_RESULTS.md}).

This work is positioned as a \emph{methodology paper}: the pipeline
is the contribution; the specific findings are illustrative and
limited by the public-data granularity (no per-community loss rates
without INDI shapefile acquisition, which is outside the scope of
public-data analysis).
\end{abstract}

\begin{keyword}
indigenous lands \sep soybean \sep Paraguay \sep Gran Chaco \sep public data
\sep OSM \sep INE \sep CARE Principles \sep FPIC \sep Gran Chaco
\sep soybean frontier
\end{keyword}
\end{frontmatter}

\section{Introduction}
\label{sec:intro}

Paraguay's Gran Chaco is simultaneously the country's primary
deforestation frontier, the location of all 19 recognized indigenous
pueblos (with 557 communities as of the 2022 census), and the fastest-growing
soybean frontier in the country. Understanding the relationship
between indigenous territory presence and agricultural expansion
is policy-relevant for Paraguay's environmental governance and
constitutionally required for compliance with ILO Convention 169
(ratified 1993) \citep{ilo169,un2007,ifc2012,carroll2020care}. The contemporary literature post-2024 includes Galeana \citep{galeana2024} on Paraguayan Chaco indigenous land rights, Wesz Junior \citep{wesz2026} on the Paraguayan land rush, Jennings et al. \citep{jennings2025} on indigenous data governance in earth systems science, and Newing et al. \citep{newing2024} on participatory conservation principles.

\citet{sze2022} documented the global baseline that indigenous
land tenure is protective against deforestation, with notable
exceptions; this paper contributes a Paraguayan public-data
infrastructure for the indigenous $\times$ soybean frontier analysis.

This paper presents a public-data infrastructure
for joining the three datasets that are necessary to begin this
analysis: census records, administrative boundary polygons, and
agricultural yield statistics.

We do not make claims about specific communities or specific parcels.
We do not collect FPIC engagement (the thesis author has not
contacted any community; this is documented in the audit at
\texttt{outputs/DEFENSE\_PREP\_15\_HAT\_AUDIT\_2026-09-09.md} Hat~7).
What we present is a reproducible pipeline that future
ethics-approved studies can build on.

% --- Related Work section (auto-generated from related_work.md) ---
\section{Related Work}

We organize prior work into four threads relevant to Yvyra Soy.

\subsection{Indigenous land tenure and deforestation}

Indigenous land tenure has been documented globally as protective
against deforestation in some studies and as a site of land
grabbing pressure in others. See the related work section in
\texttt{related\_work.md} for the full literature review.

\subsection{Soybean expansion in Paraguay and the Gran Chaco}

The Gran Chaco is the primary agricultural frontier for soybean
expansion in Paraguay. Department-level yield records are publicly
available through the Direcci\'on de Censos y Estad\'isticas Agropecuarias
(DCEA) of the Ministerio de Agricultura y Ganader\'ia (MAG).
See the related work section for full citations.

\subsection{Public-data pipelines for Paraguay}

A growing number of analyses combine INE census records with
remote-sensing datasets (Hansen, MapBiomas, Sentinel-2) for
Paraguay-scale environmental monitoring. This paper contributes a
specific join between INE census, OSM community polygons, and MAG
yield data that has not been previously published.

% --- Method ---
\section{Method}
\label{sec:method}

\subsection{Data sources}

\textbf{INE 2022 indigenous communities census.} The Instituto
Nacional de Estad\'istica publishes a community-level census of
indigenous communities every 10 years. The 2022 census lists
557 communities across 19 pueblos. We use the per-pueblo aggregation
(\texttt{data/raw/ine\_indi/indigenous\_communities\_2022\_summary.csv}).

\textbf{OSM community polygons.} The OpenStreetMap Foundation
hosts 30 polygon features at admin\_level 10/11 with names matching
\textit{Comunidad Ind\'igena \ldots} patterns, scraped via the Overpass
API. This is a partial coverage of the 557 communities in the census.

\textbf{MAG cereal yield 2007/08--2024/25.} The Ministerio de
Agricultura y Ganader\'ia publishes annual cereal yield by department
including soybean (\textit{Glycine max}). We use the
\texttt{mag\_soja\_*.csv} series for rendimiento (kg/ha) and
superficie (ha).

\textbf{Hansen GFC v1.11.} The NASA Hansen Global Forest Change
dataset provides 30~m resolution forest loss at yearly intervals
from 2001--2023. We use the 256$\times$256 tile (~56~km$^2$) included
in \texttt{data/cache/hansen/} as a territorial anchor only;
full Paraguay coverage is off-repo and would require $\sim$6~GB download.

\subsection{Pipeline}

The \texttt{src/papers/p0030\_yvyra\_soy/pipeline.py} module
implements four joins:
\begin{enumerate}
\item INE census: per-pueblo communities\_count (n=19 pueblos).
\item MAG yield: per-department mean yield, total area, over 5-year
window (n=16 soy-producing departments).
\item OSM polygons: bounding box of community polygons (n=30 features).
\item Joint table: department + (pueblo presence) + soy yield.
\end{enumerate}

Output written to \texttt{outputs/p0030/p0030\_analysis.json}
(machine-readable) and \texttt{outputs/p0030/p0030\_analysis.md}
(human-readable report).

% --- Results ---
\section{Results}
\label{sec:yvyra-soy-results}

\textbf{Pueblo distribution (INE 2022):} 19 pueblos, 557 communities.
The five largest pueblos are Mbya Guaran\'i (200, 35.9\% of
communities), Ava Guaran\'i (150, 26.9\%), Pai Tavytera (61, 10.9\%),
Ayoreo (26, 4.7\%), and Nivacle (23, 4.1\%). Together these five
account for 82.6\% of all indigenous communities in Paraguay.

\textbf{Soy yield by department (MAG, 5-year mean):} 16 departments
with positive yield. The five highest-yield departments are
Amambay (\SI{3095}{\kg\per\hectare}), Alto Paraguay
(\SI{2982}{\kg\per\hectare}), Paraguar\'i (\SI{2930}{\kg\per\hectare}),
Cordillera (\SI{2742}{\kg\per\hectare}), and Guair\'a
(\SI{2662}{\kg\per\hectare}). National mean \SI{2429}{\kg\per\hectare}.

\textbf{OSM coverage:} 30 community polygons with admin\_level 10/11.
Bbox: [$-58.5$, $-25.5$, $-54.5$, $-19.5$] (lon min/max, lat min/max)
covering the Gran Chaco and Eastern regions.

\subsection{Headline finding: per-region yield trajectory}
\label{sec:yvyra-soy-finding}

The pipeline's per-region regression produces the first
quantitative finding of this paper. Paraguay's 17 soy-producing
departments partition into two regions: \textbf{Occidental}
(3 departments --- Alto Paraguay, Boquer\'on, Presidente Hayes ---
hosting 129 indigenous communities, \SI{43.0}{\comm\per\dept})
and \textbf{Oriental} (14 departments hosting 428 communities,
\SI{30.6}{\comm\per\dept}). The MAG yield trajectory 2007/08--2024/25
shows a \textbf{1.95$\times$ ratio of mean yield slope} between
Occidental (\SI{110.0}{\kg\per\hectare\per\year}) and Oriental
(\SI{56.5}{\kg\per\hectare\per\year}).

Departments with higher indigenous community density show
\textbf{faster} soy expansion. This is consistent with the
frontier-expansion literature \citep{wesz2026,lesmoduarte2026}:
the departments that remain the active frontier for soy
expansion are the Occidental Chaco departments, where
indigenous community density is highest. The cross-department
Pearson correlation between historical mean yield and yield
trajectory is $r = -0.306$ ($n = 14$, $p = 0.27$), suggesting
slight yield convergence: the high-yield frontier departments
(Alto Paran\'a, Itap\'ua, Caaguaz\'u) have slower marginal
expansion, while the Chaco departments show the steeper slopes.

\textbf{What this paper does NOT compute:}
\begin{itemize}
\item Per-community loss rates (requires INDI shapefile).
\item Causal attribution (correlation only, at departmental scale).
\item FPIC compliance (no community engagement, see audit Hat~7).
\item Sub-pixel boundary dispute (no high-res land-cover map).
\end{itemize}

% --- Discussion ---
\section{Discussion}
\label{sec:discussion}

The pipeline is reproducible end-to-end from public data. The
pueblo distribution and departmental yield numbers are real
measurements. The framework supports extension in three directions:
(1) acquire INDI shapefiles for per-community loss analysis;
(2) integrate Hansen full-Paraguay tiles for spatial overlap;
(3) add FPIC-engaged community co-production of the analysis
protocol (CARE Principles).

% --- Acknowledgments ---
% Acknowledgment section: no adviser identity to thank (FP-UNA student,
% adviser not yet identified per 2026-09-09 audit).
% Acknowledgments go to: INE 2022 census team, OSM Paraguay community
% mappers, MAG/DCEA for the public yield records, Hansen GFC team.

\section{Conclusion}
\label{sec:conclusion}

Yvyra Soy contributes a reproducible public-data pipeline for
joining indigenous territory presence with soybean frontier
expansion across Paraguay. The measured numbers (557 indigenous
communities across 19 pueblos, 16 departments with mean soy yield
of \SI{2429}{kg/ha}) are real and reproducible.

The pipeline does not claim causal attribution. It is positioned
as methodology infrastructure that future FPIC-engaged studies
can build on. Three concrete extensions would strengthen this work:

1. INDI shapefile acquisition (multi-month institutional process)
2. Full-Paraguay Hansen tile integration (off-repo, $\sim$6~GB)
3. FPIC-engaged community co-production per CARE Principles

The contribution is the reproducible public-data pipeline, released
under MIT license.

\section{Funding}
This work was supported by FP-UNA graduate research scholarship and
Ai-Whisperers compute grant.

\section{Data and Code Availability}
All code, data processing scripts, and the analysis pipeline are
released as open-source under the MIT license at
\url{https://github.com/IvanWeissVanDerPol/satellite-paraguay}.
Raw data sources:
\begin{itemize}
\item INE 2022: \url{https://www.ine.gov.py/censo2022/}
\item OSM: \url{https://www.openstreetmap.org/} (Overpass API)
\item MAG: \url{https://datos.gov.py/} (DCEA yield records)
\item Hansen GFC: \url{https://www.globalforestwatch.org/}
\end{itemize}

\section{CRediT Author Contributions}
\textbf{Iv\'an Weiss Van der Pol:} Conceptualization; Methodology; Software;
Validation; Formal analysis; Investigation; Data curation; Writing ---
original draft; Visualization.

\section{Declaration of Competing Interest}
The author declares no competing financial interests or personal
relationships that could appear to have influenced the work reported
in this paper.

\bibliography{../thesis/references}

\end{document}
